% §2.1 Declarative Specifications and Kernel-Centric Translation
%   - First: what a spec IS — JSON contract encoding identity, file
%     partitioning, build/run/verify config. Refer to appendix for schema.
%   - Then: the KEY DESIGN CHOICE — kernel-centric translation. LLM gets
%     only translation_targets; build infra stays fixed. This is what
%     separates ParBench from ParEval-Repo (0% vs 37–76%).
%   - Then: prompt construction — structured prompt with source code,
%     target API, build command, infrastructure context. Anonymization.
%   - Finally: cross-API argument/verification handling — union patterns
%     for full-program; target spec for kernel-only (OpenCL .cl).
%
% §2.2 Augmentation Engine
%   - First: WHY — training data contamination concern (Rodinia is
%     everywhere). Need to test whether success = memorization.
%   - Then: WHAT — six AST-level transforms (list them). All backed by
%     libclang. Preserve parallel structure, change surface features.
%   - Then: HOW MUCH — five levels (L0–L4), increasing density.
%   - Finally: VALIDATION — 80/87 specs pass all levels; 87/87 at L1–L2.
%     7 failures only in omp_target at high intensity.
%
% §2.3 Harness: Build, Run, Verify
%   - First: three-stage pipeline with short-circuit on failure.
%   - Then: what each stage does (build → run → verify).
%   - Then: WHY conjunctive verification matters — example of Gaussian
%     elimination passing compile+exit but failing output check.
%   - Finally: five failure classifications → enables Section 5.2 analysis.

\section{ParBench Framework}
\label{sec:framework}

ParBench is a benchmark framework designed to evaluate LLM-based parallel code translation through automated, end-to-end correctness assessment. The framework comprises four integrated components: (1)~a declarative spec schema that encodes what constitutes correct execution for each benchmark kernel, (2)~an AST-driven augmentation engine that generates semantically-equivalent code variants to assess translation robustness in the face of code perturbations, (3)~an evaluation pipeline that orchestrates LLM-based translation and records diagnostic failure categories, and (4)~a three-stage build-run-verify pipeline that evaluates translated code. This section describes each component in detail.

\begin{figure*}[t]
\centering
\includegraphics[width=\textwidth]{parbench_architecture}
\caption{\textbf{ParBench system architecture.} API Selection (stage~0, red) identifies the target parallel APIs. Here we conducted a 35-repo survey. The Benchmark Corpus (stage~1, orange) contains the evaluation specs drawn form the benchmark suites. In this study, we curated 96 evaluation specs drawn from five suites: Rodinia (60), XSBench (4), RSBench (4), mixbench (3), and a curated HeCBench subset (25); spanning CUDA, OpenMP, OpenCL, and OpenMP target offload; 110 additional bulk HeCBench specs (not evaluated in this paper) bring the repository total to 206 specs. Each spec encodes declarative correctness contracts that configure both the Evaluation Pipeline and the Harness. Original kernel source code feeds into the Augmentation Engine (stage~2, gold), which applies six AST-driven semantics-preserving transforms at four increasing intensity levels (L1--L4). The Evaluation Pipeline (stage~3, blue) constructs a structured prompt, issues it to the LLM being evaluated, and extracts translated files from LLM output The Harness (stage~4, light blue) executes a three-stage Build-Run-Verify sequence with conjunctive correctness verification (exit code, stdout pattern matching, and, for a subset of specs, numeric or hash-based output-level checks). Results (stage~5, green) classify each task outcome.}
\label{fig:architecture}
\end{figure*}

\subsection{Spec Schema: Declarative Correctness Contracts}
\label{sec:spec-schema}

Each benchmark in ParBench is described by a JSON specification (spec) that serves as a declarative contract defining what ``correct execution'' means for that kernel--API variant. As illustrated by the spec configuration flow in Figure~\ref{fig:architecture} (stage~1$\rightarrow$3 and stage~1$\rightarrow$4), the spec provides both prompt configuration to the Evaluation Pipeline and build-run-verify configuration to the Harness. The spec encodes all information required to build, run, and verify a benchmark without manual intervention, enabling fully automated evaluation. The schema is formally defined in JSON Schema and enforced by an offline validator.

A spec is organized into the following field groups:

\textbf{Identity and provenance.} The \texttt{identity} block assigns a globally unique identifier following the convention \texttt{\{source\_suite\}-\{kernel\_name\}-\{parallel\_api\}} (e.g.,
\texttt{rodinia-bfs-cuda}) and records the target parallel API. The \texttt{provenance} block pins the exact source repository and commit hash, ensuring reproducibility across environments.

\textbf{File partitioning.} The \texttt{files} block partitions benchmark source files into three disjoint sets:
\texttt{prompt\_payload} (files provided to the LLM during translation), \texttt{support\_files} (build infrastructure such as Makefiles and shared headers, provided as read-only reference), and \texttt{verification\_only} (reference implementations never provided to the LLM). The field \texttt{translation\_targets} identifies a subset of \texttt{prompt\_payload} that the LLM must translate. This field is central to the kernel-centric translation methodology described in Section~\ref{sec:eval-pipeline}.

\textbf{Build, run, and verification.} The \texttt{build} block specifies compiler commands, environment variables, and expected executable path. The \texttt{run} block defines input configurations with arguments and timeouts. The \texttt{verification} block declares an ordered list of verification strategies evaluated in conjunction. All must pass for an overall PASS verdict (Section~\ref{sec:harness}). Five strategy types are implemented: \texttt{exit\_code}, \texttt{stdout\_pattern}, \texttt{stdout\_exclude\_pattern}, \texttt{numeric\_comparison} (scalar tolerance check against a reference), and \texttt{file\_hash} (SHA-256 digest match).

\textbf{Baseline results.} The \texttt{baseline\_results} block records the expected output from running the original implementation on a reference platform, providing ground truth for correctness evaluation.

% src: "96 specs = 60+4+4+3+25" verified via ls specs/{suite}-*.json: rodinia=60, xsbench=4, rsbench=4, mixbench=3, hecbench=135 (25 curated) (01-01 verified 2026-04-03)

We have defined 96 such specs spanning five benchmark suites: 60 from Rodinia~\cite{Rodinia2009}, 4 from XSBench~\cite{XSBench2014}, 4 from RSBench, 3 from mixbench, and 25 from a curated subset of HeCBench~\cite{HeCBench2023}. Of these, 87 pass the baseline build-run-verify pipeline; the remaining 9 (\knownfail{} specs) fail due to platform-specific issues, such as CUDA~12 API deprecations, missing system libraries, or pre-existing runtime errors, and are excluded from evaluation (Section~\ref{sec:benchmark-curation}). Listing~\ref{lst:bfs_spec_appendix} in Appendix~\ref{sec:appendix-framework} illustrates the spec structure using Rodinia's BFS kernel. 

\subsection{Evaluation Pipeline: Kernel-Centric Translation}
\label{sec:eval-pipeline}
The evaluation pipeline (Figure~\ref{fig:architecture}, blue) orchestrates LLM-based parallel code translation through a structured prompt-response-verify cycle. Its design reflects a deliberate methodological choice: \emph{kernel-centric translation}. ParBench isolates the LLM's ability to translate parallel code by feeding only the \texttt{translation\_targets} files (the parallel kernel code) to the LLM; while Makefiles, shared headers, serial baselines, and utility code remain unchanged in the target directory. This design isolates evaluation of whether an LLM can correctly map parallel programming patterns from one API to another, independent of its ability to replicate project file organization.

\textbf{Prompt structure.} Each translation prompt consists of a system message establishing the role of a parallel programming expert and a user message containing: (1)~the source kernel code, (2)~the target API and filenames to produce, (3)~the target build command, and (4)~read-only target infrastructure context so the LLM can match expected function signatures. Source support headers are included with instructions to inline definitions rather than emit unresolvable \texttt{\#include} directives. The full prompt template is in Appendix~\ref{sec:appendix-prompt}.

% FIX-3: Anonymization protocol description
%\textbf{Prompt anonymization.} To reduce the risk that benchmark-specific identifiers in the prompt trigger memorized translations, the evaluation pipeline applies six anonymization measures: (1)~the kernel name and description are omitted from the prompt entirely, (2)~source code comments are stripped from all files, (3)~source filenames are replaced with generic identifiers (``Source File 1'', etc.), (4)~target output filenames use a generic scheme (\texttt{translated\_0.ext}, \texttt{translated\_1.ext}, etc.), (5)~kernel names in build commands are replaced with generic placeholders, and (6)~support header and auxiliary files are renamed to generic identifiers. This anonymization is orthogonal to the AST-driven augmentation engine (Section~\ref{sec:augmentation-engine}), which modifies the source code \emph{structure}; prompt anonymization targets the surrounding \emph{metadata}. Full details appear in Section~\ref{sec:anonymization}.

\textbf{Cross-API argument and verification handling.} Source and target implementations may expect different command-line arguments and produce different \texttt{stdout} formats. ParBench handles this asymmetry with two branches. For \emph{full-program} translations (e.g., CUDA to OpenMP), the pipeline uses the \emph{source} spec's run arguments since LLM-translated binaries preserve the source's argument-parsing logic, and constructs a union of source and target \texttt{stdout\_pattern} regexes for verification, accommodating translations that produce either output format. For \emph{kernel-only} translations (e.g., OpenCL \texttt{.cl} kernel files are rewritten and the host binary is unchanged), the pipeline uses the \emph{target} spec's run arguments and verification patterns directly, since the host program's argument parsing and output format are unmodified. This dual handling ensures that the evaluation measures translation quality rather than penalizing correct translations for superficial format differences.

% Commenting out as it looks confusing to talk about kernel-only translation and then removing false positives. It is
% understanble by now that we translate kernel + host code.
% \textbf{Kernel-only vs.\ full-program translation.} Not all cross-API translations require the LLM to rewrite the entire program. ParBench distinguishes \emph{kernel-only} translations --- where only the device kernel files are translated while the host code remains untouched --- from \emph{full-program} translations, where the LLM rewrites all source files, including host-side API calls. The distinction arises naturally from API architecture: translations targeting OpenCL produce only \texttt{.cl} kernel files, because the OpenCL programming model separates kernel source (compiled at runtime via \texttt{clBuildProgram()}) from host orchestration code. In contrast, CUDA-to-OpenMP translations require rewriting the entire source file, replacing device kernels, memory management, and launch syntax with OpenMP pragmas and CPU-side equivalents. ParBench detects kernel-only translations automatically by checking whether all \texttt{translation\_targets} in the target spec end with \texttt{.cl}. For kernel-only translations, the target spec's native run arguments and verification patterns are used directly (since the host binary is unchanged), whereas full-program translations use the cross-API argument and verification handling described above.

% \textbf{False-positive detection for kernel-only translations.} Kernel-only translations introduce a subtle verification challenge: the host program may swallow kernel compilation errors and still produce expected output. To guard against such false positives, the verification pipeline includes a post-verification scan of \texttt{stdout} for OpenCL kernel build failure signatures. Translations that pass exit-code and stdout-pattern checks but contain evidence of kernel compilation failure are downgraded from PASS to \verifyfail{}. This specificity filter reduces false positives with zero additional false negatives, since legitimate kernel executions do not produce \texttt{clBuildProgram()} error messages.

%\textbf{Optional self-repair.} The framework supports an iterative self-repair path (dashed arrow in Figure~\ref{fig:architecture}): on failure, the pipeline can feed failure-specific diagnostics back to the LLM for a corrected translation attempt. This capability is not exercised in the evaluation reported here; each sample is a single-attempt translation with no iterative repair (Section~\ref{sec:experimental-setup}).

\textbf{Response extraction.} Each LLM response is parsed to recover the translated source files using a multi-tier extraction strategy (Appendix~\ref{sec:appendix-extraction}) that progressively relaxes matching criteria from explicit filename annotations to fuzzy and elimination-based heuristics. If any expected target file cannot be recovered or is empty, the task is classified as \extractionfail{}.

\textbf{Complexity taxonomy.} To enable stratified analysis of translation difficulty, each source--target pair is classified by cardinality of the translation: \texttt{single\_file} ($1 source {\to}1 target$), \texttt{multi\_to\_single} ($N{\to}1$), \texttt{single\_to\_multi} ($1{\to}N$), or \texttt{multi\_to\_multi} ($N{\to}M$).

\subsection{Augmentation Engine: Probing Robustness via Code Perturbations}
\label{sec:augmentation-engine}

Because widely used suites such as Rodinia~\cite{Rodinia2009} may already appear in LLM training data, the augmentation engine, the Augmentation component in Figure~\ref{fig:architecture} (gold), generates semantically-equivalent code variants by applying AST-level transformations.

The engine implements six transforms, each backed by the \texttt{libclang} API for AST analyssis to guarantee syntactic validity:
\begin{enumerate}
    \item \textbf{SwapCondition}: Reverses comparison operands in expressions (e.g., \texttt{a < b} $\to$ \texttt{b > a}), with guards against side-effect-bearing sub-expressions and assignments.
    \item \textbf{ArithmeticTransform}: Converts between compound and expanded assignments (e.g., \texttt{x += 1} $\leftrightarrow$ \texttt{x = x + 1}).
    
    \item \textbf{ChangeNames}: Renames local variables and function parameters using scope-aware reference collection, respecting OpenMP pragma clauses and boundaries.
    
    \item \textbf{TypedefExpansion}: Replaces typedef aliases with their underlying primitive types at AST-identified type-reference sites.
    \item \textbf{PointerArithmeticToArrayIndex}: Converts between pointer dereference and array indexing (e.g., \texttt{*(arr + i)} $\to$ \texttt{arr[i]}).
    
    \item \textbf{ChangeFunctionNames}: Renames \texttt{static} helper functions across definitions and all call sites within the translation unit: \texttt{\_\_global\_\_}/\texttt{\_\_kernel\_\_} entry points are excluded to avoid breaking OpenCL runtime name resolution.
\end{enumerate}

These six transforms target surface-level code features: naming conventions, syntax sugar, type aliases, and expression forms, that are most likely to be memorized verbatim by LLMs trained on open-source repositories, while preserving the parallel structure (thread indexing, synchronization, memory management) that is the subject of evaluation. Parallelism-aware transforms such as loop reordering or reduction-strategy substitution are out of scope because they risk altering the observable output and would require per-kernel correctness proofs.

All transforms operate on the parsed AST rather than raw text, apply to every source file in the prompt payload (including both host code and GPU kernel function bodies, excluding \texttt{\#include} headers), and each candidate rewrite is validated by re-parsing the transformed code to confirm no new errors are introduced.

% [Float moved to Appendix D]
%The complete augmentation level definitions are provided in Appendix Table~\ref{tab:augmentation-levels}.
Five augmentation levels (L0--L4) control the density of applied transforms. At L0, the source code is unmodified. At L1, exactly one transform is selected and applied to a single candidate site. Levels L2--L4 increase both the number of active transforms and the fraction of eligible sites rewritten, controlled by a density parameter $f$ ($f{=}0.33$ at L2, $f{=}0.66$ at L3, $f{=}1.0$ at L4). Within each selected transform, the same fraction determines how many eligible AST nodes to rewrite: $\max(1, \lfloor n \cdot f \rfloor)$ sites are modified, where $n$ is the number of eligible candidate sites for that transform in the file. The complete level definitions are in Appendix~\ref{sec:appendix-framework}.

% src: Rodinia 53/53 non-KF pass L1-L4 (retest_post_session2.json, 2026-03-20).
% src: Non-Rodinia 27/34 pass L1-L4; 34/34 pass L1-L2 (non_rodinia_sweep.json, 2026-04-26).
% Combined: 80/87 pass ALL L1-L4; 87/87 pass L1-L2. 7 failures = omp_target at L3-L4 only.

The augmentation baseline validates semantics preservation: 80 of 87 non-\knownfail{} specs across all five suites pass the build-run-verify pipeline at every level L1--L4, with zero correctness regressions relative to L0. All 87 pass at L1--L2. The 7 high-intensity failures occur exclusively in \texttt{omp\_target} variants at L3 or L4, where aggressive variable renaming interferes with OpenMP target data-mapping clauses. This is a known limitation of the current transform guards.
%Transform application frequency across 240 augmented tasks: SwapCondition (162), ArithmeticTransform (69), ChangeNames (55), TypedefExpansion (7), PointerArithmeticToArrayIndex (6), ChangeFunctionNames (2).

\subsection{Harness Pipeline: Build, Run, Verify}
\label{sec:harness}

The harness pipeline is a three-stage evaluation corresponding to the Build, Run, and Verify sub-stages within the Harness Pipeline component of Figure~\ref{fig:architecture} (light blue). Each stage consumes the output of the previous stage and failure at any stage short-circuits subsequent stages. The Harness Pipeline receives translated files from the Evaluation Pipeline and build-run-verify configuration from the Spec JSON, and then emits a verdict to the Result node.

\textbf{Build stage.} The builder resolves the spec's working directory, executes an optional \texttt{clean} command, substitutes template variables (e.g., \texttt{\$\{CUDA\_DIR\}}) from a platform-specific configuration file that maps tokens to local compiler and SDK paths, and invokes the build command. After compilation, the builder verifies that the expected executable exists on the disk. The build stage captures both \texttt{stdout} and \texttt{stderr} for diagnostic reporting and enforces a configurable timeout (default: 600 seconds).

\textbf{Run stage.} The runner executes the compiled binary with arguments drawn from the spec's named input configuration (typically \texttt{correctness}). Execution proceeds without a shell to avoid argument-parsing ambiguities. The runner captures \texttt{stdout}, \texttt{stderr}, \texttt{exit code}, and wall-clock duration, and enforces the spec-defined timeout (default: 300 seconds). An optional CPU-time measurement mode captures user+system time independently of wall-clock duration.

\textbf{Verify stage.} The verifier applies the spec's verification strategies in declared order. For \texttt{exit\_code} verification, the actual \texttt{exit code} is compared against the expected value. For \texttt{stdout\_pattern} verification, a regular expression is applied to the captured \texttt{stdout}. A FAIL or ERROR result triggers immediate termination; a PASS is recorded and evaluation continues to the next strategy. All implemented strategies must pass for the overall verdict to be PASS. In cross-API translation evaluation, the verification stage uses a combined pattern that accepts output matching either the source or target spec's \texttt{stdout} format (Section~\ref{sec:eval-pipeline}), preventing false negatives from format differences between API implementations. 

The conjunction of \texttt{exit code} and \texttt{stdout}-pattern verification is essential for accurate measurement. For example, a CUDA translation of the Gaussian elimination kernel (from OpenCL) compiles without error, runs to completion with \texttt{exit code} zero, and produces partial \texttt{stdout} output, but fails to print the expected \texttt{Total:} timing summary that the reference implementation produces. This indicates a systematic translation defect rather than a transient failure. Compilation-only verification~\cite{CodeRosetta2024} would misclassify such a task as a successful translation; exit-code-only verification would similarly fail, since this translation terminated normally without signaling an error.

For kernels that produce quantitative output, stronger verification is available: \texttt{numeric\_comparison} extracts a scalar from stdout via a configurable regex and compares it against a reference value within a per-spec tolerance, while \texttt{file\_hash} verifies that a designated output file matches a reference SHA-256 digest. Of the 87 non-\knownfail{} specs, 7 use these stronger oracle; the remaining 80 rely on exit-code and stdout-pattern checks that capture each benchmark's own self-checking output (e.g., printed \texttt{PASS}/\texttt{FAIL} verdicts or expected timing summaries). Because the dominant failure mode is build-stage API adaptation (Section~\ref{sec:results}), the current oracle distribution is sufficient to classify the majority of outcomes; upgrading weak oracles to numeric verification for remaining specs is identified as future work (Section~\ref{sec:discussion}).

Each harness invocation produces one of four failure classifications: BUILD\_FAIL, RUN\_FAIL, VERIFY\_FAIL, or TIMEOUT, providing diagnostic granularity that enables systematic failure-mode analysis. The evaluation pipeline (Section~\ref{sec:eval-pipeline}) adds a fifth classification, EXTRACTION\_FAIL, when the LLM response does not contain parseable target files.

% A \texttt{--resume} flag allows batch evaluations to skip previously completed tasks, supporting append-only result accumulation without re-executing passing benchmarks.



{\color{blue}
\section{ParBench Framework}
\label{sec:framework}
To better evaluate LLM-based translation of parallel programming APIs, ParBench defines a controlled, kernel-centric setting that separates parallel kernel translation from repository-level challenges such as build-system reconstruction. The framework realizes this setting through four components: (1)~declarative specifications, (2)~a build-run-verify harness, (3)~an AST-driven augmentation engine, and (4)~an evaluation pipeline that orchestrates LLM translation and records diagnostic failure categories. Together, these components enable direct measurement of parallel API translation capability under end-to-end executability constraints while supporting controlled robustness analysis. Figure~\ref{fig:architecture} shows the overall architecture.

% ParBench evaluates LLM-based parallel code translation through automated, end-to-end correctness assessment. The framework has four components: (1)~declarative specifications that define correct execution for each kernel, (2)~a build-run-verify harness, (3)~an AST-driven augmentation engine for robustness testing, and (4)~an evaluation pipeline that orchestrates LLM translation and records failure categories. Figure~\ref{fig:architecture} shows the overall architecture.

\begin{figure*}[t]
\centering
\includegraphics[width=\textwidth]{parbench_architecture}
\caption{\textbf{ParBench architecture.} The pipeline proceeds from API selection (stage~0) through a curated benchmark corpus (stage~1, 96~specs from five suites), optional AST-driven augmentation (stage~2), LLM-based translation (stage~3), and a build-run-verify harness (stage~4) that classifies each outcome. The dashed arrow indicates an optional self-repair path not exercised in this evaluation.}
\label{fig:architecture}
\end{figure*}

\subsection{Declarative Specifications and Kernel-Centric Translation}
\label{sec:spec-schema}

Each benchmark instance in ParBench is encoded as a JSON specification that serves as a declarative contract for building, running, and verifying a kernel--API variant. A spec declares the kernel's identity and provenance, including the pinned repository and commit; partitions source files into \texttt{translation\_targets} (files the LLM must translate), \texttt{support\_files} (read-only build and interface context), and \texttt{verification\_only} (reference implementations never shown to the LLM); and defines the build command, run arguments, and verification strategies. The schema is formally validated offline using JSON Schema draft-07. An example spec and the full structure are provided in Appendix~\ref{sec:appendix-framework-schema}.

The key design choice is \textbf{kernel-centric translation}: the LLM is asked to modify only the \texttt{translation\_targets} files and produce functionally equivalent target-API code, while Makefiles, host scaffolding, and verification logic remain fixed. This isolates parallel API translation from repository reconstruction---the confound that drives ParEval-Repo's 0\% pass rate on applications above 133~SLoC~\cite{ParEvalRepo2025}. The evaluation pipeline constructs a structured prompt containing the source kernel code, the target API and output filenames, the build command, and read-only target infrastructure context, with prompt anonymization to reduce benchmark identity leakage (Appendix~\ref{sec:appendix-prompt}). For cross-API translations, the pipeline handles argument and verification asymmetry automatically: full-program translations use a union of source and target verification patterns, whereas kernel-only translations, such as OpenCL \texttt{.cl} files, use the target spec directly.

\subsection{AST-driven Augmentation Engine}
\label{sec:augmentation-engine}

Widely used suites such as Rodinia~\cite{Rodinia2009} are likely present in LLM training data, so success on unmodified source may reflect memorization. The augmentation engine generates semantically equivalent but syntactically novel variants by applying six AST-level transforms backed by the libclang API: condition operand swapping, arithmetic form conversion, scope-aware variable renaming, typedef expansion, pointer-to-array notation conversion, and static function renaming. All transforms preserve parallel structure (thread indexing, synchronization, memory management) while changing surface-level code features most likely to be memorized verbatim. Parallelism-aware transforms (e.g., loop reordering) are excluded because they risk altering observable output.

Five augmentation levels (L0--L4) control density. L0 is unmodified source. L1 applies one transform to one site. L2--L4 apply increasing fractions of transforms and candidate sites ($f{=}0.33$, $0.66$, $1.0$). Baseline validation confirms semantic preservation: 80 of 87 non-\knownfail{} specs pass at all levels L1--L4, with zero correctness regressions. All 87 pass at L1--L2. The 7 high-intensity failures occur exclusively in \texttt{omp\_target} variants where aggressive variable renaming interferes with OpenMP target data-mapping clauses. Level definitions and per-transform details are in Appendix Tables~\ref{tab:augmentation-levels} and~\ref{sec:appendix-c3}.

\subsection{Harness: Build, Run, Verify}
\label{sec:harness}

The harness is a three-stage pipeline where failure at any stage short-circuits subsequent stages. The \textbf{build stage} resolves the spec's working directory, substitutes platform-specific paths, and invokes compilation. The \textbf{run stage} executes the compiled binary with spec-defined arguments and timeout. The \textbf{verify stage} applies an ordered list of verification strategies in conjunction---all must pass for a \pass{} verdict. Five strategy types are implemented: \texttt{exit\_code}, \texttt{stdout\_pattern}, \texttt{stdout\_exclude\_pattern}, \texttt{numeric\_comparison}, and \texttt{file\_hash}.

Conjunctive verification is essential. For example, a CUDA translation of Gaussian elimination compiles, runs to completion with exit code zero, and produces partial stdout---but omits the expected timing summary. Compile-only~\cite{CodeRosetta2024} or exit-code-only verification would misclassify this as successful. Each outcome receives one of five failure classifications---\buildfail{}, \runfail{}, \verifyfail{}, \extractionfail{}, or \pass{}---enabling systematic failure-mode analysis (Section~\ref{sec:failure-taxonomy}).

}